\volpart{第二编\quad 响应、无序与关联电子（8.512）}
\originalchapter{因果响应、涨落与散射}
\label{ch:response}
\source{8.512课程页第1--4讲；Mark Spencer Rudner讲义及Lee讲义。}
\originalsection{从微扰演化到Kubo公式}
设未扰动密度矩阵为$\rho_0=\e^{-\beta H_0}/Z$，施加$H'(t)=-f(t)B$。在相互作用绘景，一阶演化算符为
\begin{equation}
U_I(t)\simeq1+\frac\ii\hbar\int_{-\infty}^t f(t')B_I(t')\dd t'.
\end{equation}
将$U_I^\dagger A_IU_I$展开，响应为
\begin{equation}
\delta\avg{A(t)}=\int\dd t'\,\chi_{AB}(t-t')f(t'),\qquad\chi_{AB}(t)=\frac\ii\hbar\theta(t)\avg{[A(t),B(0)]}.
\end{equation}
$\theta(t)$表达原因不能晚于结果。此处采用对$-fB$响应的正静态磁化率约定。OCW若以$+fB$扰动定义响应函数，则其$\chi$为这里的负号，尤其题面中的$S=-2(n_B+1)\chi''$不能与本章符号无说明混用。Fourier约定$\chi(\omega)=\int\e^{\ii\omega t}\chi(t)\dd t$。

插入$H_0\ket n=E_n\ket n$，对$A=B^\dagger$得谱表示
\begin{equation}
\chi_{AA^\dagger}(\omega)=\sum_{nm}\frac{(p_n-p_m)|A_{nm}|^2}{E_m-E_n-\hbar\omega-\ii0^+},\qquad p_n=\frac{\e^{-\beta E_n}}Z.
\end{equation}
分母极点位于复频率下半面，所以$\chi$在上半面解析。对实扰动与Hermitian观测量，$\chi(-\omega)=\chi(\omega)^*$。复平面积分给Kramers--Kronig关系$\chi'(\omega)=\pi^{-1}\mathcal P\int\chi''(\omega')/(\omega'-\omega)\dd\omega'$，以及实虚部互换的负号关系。关系要求响应在无穷远充分衰减；有常数瞬时项时应先扣去该项。
\originalsection{涨落耗散定理}
定义非对称谱$S_{AA^\dagger}(\omega)=\int\e^{\ii\omega t}\avg{A(t)A^\dagger(0)}\dd t$。插入能量态可得$S(\omega)=2\pi\hbar\sum_{nm}p_n|A_{nm}|^2\delta(E_m-E_n-\hbar\omega)$。交换$n,m$后，$p_m=\e^{-\beta\hbar\omega}p_n$，所以$S_{A^\dagger A}(-\omega)=\e^{-\beta\hbar\omega}S_{AA^\dagger}(\omega)$。这正是详细平衡：零温体系能吸收外界能量，却不能从基态自发释放正能量。

利用谱分母的虚部$\pi\delta(E_m-E_n-\hbar\omega)$，得到
\begin{equation}
S(\omega)=\frac{2\hbar}{1-\e^{-\beta\hbar\omega}}\chi''(\omega),\qquad S_{\rm sym}(\omega)=\hbar\coth\frac{\beta\hbar\omega}2\,\chi''(\omega).
\end{equation}
非对称谱对应实验能量吸收与发射，不能无条件换成对称谱。经典低频极限$S_{\rm sym}\simeq2\kb T\chi''/\omega$，热涨落与耗散系数相联。系统必须处于热平衡；在驱动稳态只测到噪声并不能自动套此式。
\originalsection{密度相关与中子散射}
中子通过核的短程散射长度$b$与晶格密度$\rho_Q=\sum_j\e^{-\ii Q\cdot r_j}$耦合。Born近似把散射率写成常数乘$S(Q,\omega)$，其中$Q=k_i-k_f$、$\hbar\omega$是中子失去的能量。结构因子由$\avg{\rho_Q(t)\rho_{-Q}(0)}$给出，既含弹性Bragg峰也含声子等集体激发。

谐振位移满足高斯统计，但不同时间的位移不对易。令$X=-\ii Q u_j(t)$、$Y=\ii Q u_l(0)$，由于$[X,Y]$是数，Baker--Campbell--Hausdorff公式为$\e^X\e^Y=\e^{X+Y}\e^{[X,Y]/2}$。这项量子相位保证声子产生权重$n_q+1$与消灭权重$n_q$不同。展开位移指数到最低非弹性阶得
\begin{align}
S(Q,\omega)\propto\e^{-2W}\bigg[&\sum_G\delta(Q-G)\delta(\omega)\\
+&\sum_{q\lambda}\frac{\hbar|Q\cdot e_{q\lambda}|^2}{2M\omega_{q\lambda}}\big\{(n_q+1)\sum_G\delta(Q-G-q)\delta(\omega-\omega_q)\nonumber\\
&\hspace{35mm}+n_q\sum_G\delta(Q-G+q)\delta(\omega+\omega_q)\big\}\bigg].\nonumber
\end{align}
第一项是弹性峰，后二项是产生与吸收一个声子。选不同方向$Q$可测极化$e_{q\lambda}$，因为强度含投影平方，不能只凭频率判定纵模或横模。
\originalsection{双对易与频率和规则}
对能量本征态$\ket0$与Hermitian算符$A$，插入完备关系得到$\sum_n(E_n-E_0)|A_{n0}|^2=\avg{[A,[H,A]]}/2$。取$A=x$、$H=p^2/(2m)+V(x)$，内对易$[H,x]=-\ii\hbar p/m$，双对易为$\hbar^2/m$，因此$\sum_n(E_n-E_0)|x_{n0}|^2=\hbar^2/(2m)$。这不依赖具体势能。

对密度$\rho_q=\sum_j\e^{-\ii q\cdot r_j}$，位置依赖相互作用与密度对易，动能给$\avg{[\rho_{-q},[H,\rho_q]]}=N\hbar^2q^2/m$。于是密度谱的一阶频率矩由粒子数与裸质量固定，电子相互作用只能重新分配谱权重。连续电子系统的光学和规则$\int_0^\infty\operatorname{Re}\sigma(\omega)\dd\omega=\pi ne^2/(2m)$还包含零频$\delta$项的半轴权重。只截取一条晶格能带时，右侧改由该带动能曲率给出，不能将裸质量规则直接当作有限带模型内部的精确恒等式。
